% The debug report of spec section 17.3. Build with `make reports` from the
% repository root, `make` in this directory, or
% `latexmk -pdf -cd report_debug/debug_report.tex`.
%
% Grouped by theme rather than by date, because chronology is how the problems
% were met and theme is how they are useful to anybody else. The source of every
% entry is docs/ENGINEERING_LOG.md, which was written as the work happened.
%
% The preamble and the generated value macros are shared with the main report.
% See ../report/preamble.tex for the dash rule and the number rule that govern
% every LaTeX source in this repository.
\documentclass[11pt,a4paper,oneside]{report}

% Shared preamble for all three reports of spec section 17.
%
% Two mechanical rules govern every LaTeX source in this repository, and both
% are enforced in CI rather than by review. They are documented here once,
% because a reader of a chapter will otherwise wonder why the source looks the
% way it does.
%
% THE DASH RULE. Ground rule 13 forbids U+2014 and U+2013 anywhere in the
% project, and forbids the LaTeX ligature forms in .tex sources as well, since
% those typeset to exactly the characters the rule bans.
% scripts/check_dashes.py rejects any run of two or more adjacent hyphens in a
% .tex file. That collides with the one place this project genuinely needs two
% hyphens on the page: a command line flag. The collision is resolved in
% source rather than by an exemption. The macro below inserts an empty group
% between the two hyphens, so the source never carries the adjacent pair and
% the typeset page carries exactly two hyphen glyphs:
%
%     \texttt{\dd{}engine rules}     typesets as the flag a reader would type
%
% Inside a verbatim block the same macro is available because the shell
% environment declares the backslash and braces as command characters.
%
% THE NUMBER RULE. Law 3 forbids a hand typed number anywhere in a report.
% scripts/check_traceability.py scans these sources for anything shaped like a
% measurement and fails unless the line names the committed result file the
% number came from. The chapters of this report therefore contain no digits at
% all: every quoted number is a macro defined in tables/values.tex by
% scripts/make_report_tables.py, beside the path it was read from, and every
% table and figure is generated by that script or by
% scripts/make_report_figures.py. Editing a number by hand is not a policy
% violation that somebody might notice. It is a build failure.

\usepackage[T1]{fontenc}
\usepackage[utf8]{inputenc}
\usepackage{lmodern}
\usepackage[margin=2.5cm]{geometry}
\usepackage{microtype}
\usepackage{booktabs}
\usepackage{longtable}
\usepackage{array}
\usepackage{graphicx}
\usepackage{xcolor}
\usepackage[font=small,labelfont=bf]{caption}
\usepackage{enumitem}
\usepackage{fancyvrb}
\usepackage{parskip}
\usepackage[hidelinks]{hyperref}

% Result paths are printed verbatim and are allowed to break at a slash. The
% url package reads its argument with its own catcodes, so an underscore in a
% run id needs no escape, which matters: an escaped path in the source would no
% longer match the citation pattern the traceability check looks for, and the
% check would silently stop seeing the citation it was given.
%
% Run identifiers are long and carry no slash, so the break list is widened
% inside the command's own style: a run id that cannot break would overflow its
% column rather than wrapping, and the paths are the one thing in these
% documents that a reader may need to type back in.
\DeclareUrlCommand\rpath{%
  \urlstyle{tt}%
  \def\UrlBreaks{\do\/\do\.\do\-\do\_\do\T\do\Z\do\p\do\v\do\b\do\a}%
  \def\UrlBigBreaks{\do\/}%
}

% An identifier printed verbatim in a narrow table column: a metric name, a
% slice name, a taxonomy label. Same mechanism as \rpath, with breaks allowed
% at the separators identifiers actually use, so a long label wraps inside its
% column instead of running off the page.
\DeclareUrlCommand\ident{%
  \urlstyle{tt}%
  \def\UrlBreaks{\do\_\do\-\do\/\do\:\do\.}%
  \def\UrlBigBreaks{}%
}

% The two hyphen macro described at the top of this file.
\newcommand{\dd}{-{}-}

% A command line flag, written as the reader would type it.
\newcommand{\flag}[1]{\texttt{\dd{}#1}}

% Console transcripts. commandchars lets \dd{} through so that a flag can be
% shown inside an otherwise verbatim block.
\DefineVerbatimEnvironment{shell}{Verbatim}{%
  commandchars=\\\{\},%
  fontsize=\small,%
  frame=single,%
  framesep=3mm,%
  rulecolor=\color{black!30}%
}

% Markup samples, with no command characters at all, so angle brackets, quotes
% and braces are exactly what they look like.
\DefineVerbatimEnvironment{markup}{Verbatim}{%
  fontsize=\small,%
  frame=single,%
  framesep=3mm,%
  rulecolor=\color{black!30}%
}

% A finding code, a taxonomy label, or any other identifier that has one exact
% spelling and must not be hyphenated across a line break.
\newcommand{\code}[1]{\texttt{#1}}

% The standing caveat, printed wherever a reader could otherwise mistake a
% measurement on this corpus for a claim about the open web.
\newcommand{\syntheticnote}{%
  \emph{Measured on a seeded synthetic corpus generated by this repository.
  Behaviour on real world pages is unmeasured and is stated as unmeasured.}%
}

\setlist{itemsep=2pt,parsep=0pt,topsep=4pt}
\captionsetup[table]{skip=4pt}
\captionsetup[figure]{skip=6pt}

% Generated by scripts/make_report_tables.py. Do not edit.
% Every definition carries the committed file its value was read from.
\newcommand{\vRulesMacroF}{0.7911}  % results: experiments/results/test/2026-08-26T21-15-35Z_p6-rules_6457ac7/metrics.json
\newcommand{\vRulesMicroF}{0.7678}  % results: experiments/results/test/2026-08-26T21-15-35Z_p6-rules_6457ac7/metrics.json
\newcommand{\vRulesSeenMacroF}{0.7991}  % results: experiments/results/test/2026-08-26T21-15-35Z_p6-rules_6457ac7/metrics.json
\newcommand{\vRulesUnseenMacroF}{0.7276}  % results: experiments/results/test/2026-08-26T21-15-35Z_p6-rules_6457ac7/metrics.json
\newcommand{\vRulesAbstention}{0.2620}  % results: experiments/results/test/2026-08-26T21-15-35Z_p6-rules_6457ac7/metrics.json
\newcommand{\vRulesCommittedAccuracy}{0.9702}  % results: experiments/results/test/2026-08-26T21-15-35Z_p6-rules_6457ac7/metrics.json
\newcommand{\vRulesLatencyFifty}{35.2}  % results: experiments/results/test/2026-08-26T21-15-35Z_p6-rules_6457ac7/metrics.json
\newcommand{\vRulesLatencyNinetyFive}{134.0}  % results: experiments/results/test/2026-08-26T21-15-35Z_p6-rules_6457ac7/metrics.json
\newcommand{\vRulesLoadShare}{0.9912}  % results: experiments/results/test/2026-08-26T21-15-35Z_p6-rules_6457ac7/metrics.json
\newcommand{\vRulesTotalSeconds}{137.4}  % results: experiments/results/test/2026-08-26T21-15-35Z_p6-rules_6457ac7/metrics.json
\newcommand{\vRulesCleanMacroF}{0.8677}  % results: experiments/results/test/2026-08-26T21-15-35Z_p6-rules_6457ac7/metrics.json
\newcommand{\vRulesPartialMacroF}{0.8677}  % results: experiments/results/test/2026-08-26T21-15-35Z_p6-rules_6457ac7/metrics.json
\newcommand{\vRulesMixedMacroF}{0.8005}  % results: experiments/results/test/2026-08-26T21-15-35Z_p6-rules_6457ac7/metrics.json
\newcommand{\vRulesHostileMacroF}{0.4681}  % results: experiments/results/test/2026-08-26T21-15-35Z_p6-rules_6457ac7/metrics.json
\newcommand{\vRulesMissingPrecision}{1.0000}  % results: experiments/results/test/2026-08-26T21-15-35Z_p6-rules_6457ac7/metrics.json
\newcommand{\vRulesMissingRecall}{0.5684}  % results: experiments/results/test/2026-08-26T21-15-35Z_p6-rules_6457ac7/metrics.json
\newcommand{\vRulesMissingAccusations}{449}  % results: experiments/results/test/2026-08-26T21-15-35Z_p6-rules_6457ac7/metrics.json
\newcommand{\vRulesMissingCorrect}{449}  % results: experiments/results/test/2026-08-26T21-15-35Z_p6-rules_6457ac7/metrics.json
\newcommand{\vRulesMissingWrong}{0}  % results: experiments/results/test/2026-08-26T21-15-35Z_p6-rules_6457ac7/metrics.json
\newcommand{\vRulesWrongAccusations}{28}  % results: experiments/results/test/2026-08-26T21-15-35Z_p6-rules_6457ac7/metrics.json
\newcommand{\vRulesWrongRecall}{0.9333}  % results: experiments/results/test/2026-08-26T21-15-35Z_p6-rules_6457ac7/metrics.json
\newcommand{\vNgramMacroF}{0.7394}  % results: experiments/results/test/2026-08-26T21-17-53Z_p6-ngram_6457ac7/metrics.json
\newcommand{\vNgramMicroF}{0.8226}  % results: experiments/results/test/2026-08-26T21-17-53Z_p6-ngram_6457ac7/metrics.json
\newcommand{\vNgramSeenMacroF}{0.7735}  % results: experiments/results/test/2026-08-26T21-17-53Z_p6-ngram_6457ac7/metrics.json
\newcommand{\vNgramUnseenMacroF}{0.5637}  % results: experiments/results/test/2026-08-26T21-17-53Z_p6-ngram_6457ac7/metrics.json
\newcommand{\vNgramAbstention}{0.0363}  % results: experiments/results/test/2026-08-26T21-17-53Z_p6-ngram_6457ac7/metrics.json
\newcommand{\vNgramCommittedAccuracy}{0.8204}  % results: experiments/results/test/2026-08-26T21-17-53Z_p6-ngram_6457ac7/metrics.json
\newcommand{\vNgramLatencyFifty}{165.6}  % results: experiments/results/test/2026-08-26T21-17-53Z_p6-ngram_6457ac7/metrics.json
\newcommand{\vNgramLatencyNinetyFive}{261.5}  % results: experiments/results/test/2026-08-26T21-17-53Z_p6-ngram_6457ac7/metrics.json
\newcommand{\vNgramLoadShare}{0.9874}  % results: experiments/results/test/2026-08-26T21-17-53Z_p6-ngram_6457ac7/metrics.json
\newcommand{\vNgramTotalSeconds}{137.9}  % results: experiments/results/test/2026-08-26T21-17-53Z_p6-ngram_6457ac7/metrics.json
\newcommand{\vNgramCleanMacroF}{0.9350}  % results: experiments/results/test/2026-08-26T21-17-53Z_p6-ngram_6457ac7/metrics.json
\newcommand{\vNgramPartialMacroF}{0.9290}  % results: experiments/results/test/2026-08-26T21-17-53Z_p6-ngram_6457ac7/metrics.json
\newcommand{\vNgramMixedMacroF}{0.7151}  % results: experiments/results/test/2026-08-26T21-17-53Z_p6-ngram_6457ac7/metrics.json
\newcommand{\vNgramHostileMacroF}{0.5409}  % results: experiments/results/test/2026-08-26T21-17-53Z_p6-ngram_6457ac7/metrics.json
\newcommand{\vNgramMissingPrecision}{0.9731}  % results: experiments/results/test/2026-08-26T21-17-53Z_p6-ngram_6457ac7/metrics.json
\newcommand{\vNgramMissingRecall}{0.4114}  % results: experiments/results/test/2026-08-26T21-17-53Z_p6-ngram_6457ac7/metrics.json
\newcommand{\vNgramMissingAccusations}{334}  % results: experiments/results/test/2026-08-26T21-17-53Z_p6-ngram_6457ac7/metrics.json
\newcommand{\vNgramMissingCorrect}{325}  % results: experiments/results/test/2026-08-26T21-17-53Z_p6-ngram_6457ac7/metrics.json
\newcommand{\vNgramMissingWrong}{9}  % results: experiments/results/test/2026-08-26T21-17-53Z_p6-ngram_6457ac7/metrics.json
\newcommand{\vNgramWrongAccusations}{23}  % results: experiments/results/test/2026-08-26T21-17-53Z_p6-ngram_6457ac7/metrics.json
\newcommand{\vNgramWrongRecall}{0.7000}  % results: experiments/results/test/2026-08-26T21-17-53Z_p6-ngram_6457ac7/metrics.json
\newcommand{\vLlmMacroF}{0.6353}  % results: experiments/results/test/2026-08-26T20-32-48Z_p6-llm_6457ac7/metrics.json
\newcommand{\vLlmMicroF}{0.6953}  % results: experiments/results/test/2026-08-26T20-32-48Z_p6-llm_6457ac7/metrics.json
\newcommand{\vLlmSeenMacroF}{0.6373}  % results: experiments/results/test/2026-08-26T20-32-48Z_p6-llm_6457ac7/metrics.json
\newcommand{\vLlmUnseenMacroF}{0.6045}  % results: experiments/results/test/2026-08-26T20-32-48Z_p6-llm_6457ac7/metrics.json
\newcommand{\vLlmAbstention}{0.0255}  % results: experiments/results/test/2026-08-26T20-32-48Z_p6-llm_6457ac7/metrics.json
\newcommand{\vLlmCommittedAccuracy}{0.7037}  % results: experiments/results/test/2026-08-26T20-32-48Z_p6-llm_6457ac7/metrics.json
\newcommand{\vLlmLatencyFifty}{524189.3}  % results: experiments/results/test/2026-08-26T20-32-48Z_p6-llm_6457ac7/metrics.json
\newcommand{\vLlmLatencyNinetyFive}{674727.4}  % results: experiments/results/test/2026-08-26T20-32-48Z_p6-llm_6457ac7/metrics.json
\newcommand{\vLlmLoadShare}{0.0537}  % results: experiments/results/test/2026-08-26T20-32-48Z_p6-llm_6457ac7/metrics.json
\newcommand{\vLlmTotalSeconds}{2565.8}  % results: experiments/results/test/2026-08-26T20-32-48Z_p6-llm_6457ac7/metrics.json
\newcommand{\vLlmCleanMacroF}{0.6696}  % results: experiments/results/test/2026-08-26T20-32-48Z_p6-llm_6457ac7/metrics.json
\newcommand{\vLlmPartialMacroF}{0.6522}  % results: experiments/results/test/2026-08-26T20-32-48Z_p6-llm_6457ac7/metrics.json
\newcommand{\vLlmMixedMacroF}{0.6352}  % results: experiments/results/test/2026-08-26T20-32-48Z_p6-llm_6457ac7/metrics.json
\newcommand{\vLlmHostileMacroF}{0.5265}  % results: experiments/results/test/2026-08-26T20-32-48Z_p6-llm_6457ac7/metrics.json
\newcommand{\vLlmMissingPrecision}{0.7803}  % results: experiments/results/test/2026-08-26T20-32-48Z_p6-llm_6457ac7/metrics.json
\newcommand{\vLlmMissingRecall}{0.6835}  % results: experiments/results/test/2026-08-26T20-32-48Z_p6-llm_6457ac7/metrics.json
\newcommand{\vLlmMissingAccusations}{692}  % results: experiments/results/test/2026-08-26T20-32-48Z_p6-llm_6457ac7/metrics.json
\newcommand{\vLlmMissingCorrect}{540}  % results: experiments/results/test/2026-08-26T20-32-48Z_p6-llm_6457ac7/metrics.json
\newcommand{\vLlmMissingWrong}{152}  % results: experiments/results/test/2026-08-26T20-32-48Z_p6-llm_6457ac7/metrics.json
\newcommand{\vLlmWrongAccusations}{125}  % results: experiments/results/test/2026-08-26T20-32-48Z_p6-llm_6457ac7/metrics.json
\newcommand{\vLlmWrongRecall}{0.9667}  % results: experiments/results/test/2026-08-26T20-32-48Z_p6-llm_6457ac7/metrics.json
\newcommand{\vTestFields}{2317}  % results: experiments/results/bench/2026-08-26T20-32-48Z_bench_6457ac7/benchmark.json
\newcommand{\vTestForms}{240}  % results: experiments/results/bench/2026-08-26T20-32-48Z_bench_6457ac7/benchmark.json
\newcommand{\vTestTemplates}{10}  % results: experiments/results/bench/2026-08-26T20-32-48Z_bench_6457ac7/benchmark.json
\newcommand{\vLabelsObserved}{36}  % results: experiments/results/bench/2026-08-26T20-32-48Z_bench_6457ac7/benchmark.json
\newcommand{\vMissingNeeded}{790}  % results: experiments/results/bench/2026-08-26T20-32-48Z_bench_6457ac7/benchmark.json
\newcommand{\vWrongNeeded}{30}  % results: experiments/results/bench/2026-08-26T20-32-48Z_bench_6457ac7/benchmark.json
\newcommand{\vLlmWrongPrecision}{0.2320}  % results: experiments/results/bench/2026-08-26T20-32-48Z_bench_6457ac7/benchmark.json
\newcommand{\vReportingMinimum}{30}  % results: experiments/results/bench/2026-08-26T20-32-48Z_bench_6457ac7/benchmark.json
\newcommand{\vLlmRequests}{480}  % results: experiments/results/bench/2026-08-26T20-32-48Z_bench_6457ac7/benchmark.json
\newcommand{\vLlmAttempts}{480}  % results: experiments/results/bench/2026-08-26T20-32-48Z_bench_6457ac7/benchmark.json
\newcommand{\vLlmRetried}{0}  % results: experiments/results/bench/2026-08-26T20-32-48Z_bench_6457ac7/benchmark.json
\newcommand{\vLlmFailed}{0}  % results: experiments/results/bench/2026-08-26T20-32-48Z_bench_6457ac7/benchmark.json
\newcommand{\vLlmUnresolved}{0}  % results: experiments/results/bench/2026-08-26T20-32-48Z_bench_6457ac7/benchmark.json
\newcommand{\vLlmPromptTokens}{677632}  % results: experiments/results/bench/2026-08-26T20-32-48Z_bench_6457ac7/benchmark.json
\newcommand{\vLlmCompletionTokens}{159126}  % results: experiments/results/bench/2026-08-26T20-32-48Z_bench_6457ac7/benchmark.json
\newcommand{\vLlmMinutes}{42.8}  % results: experiments/results/bench/2026-08-26T20-32-48Z_bench_6457ac7/benchmark.json
\newcommand{\vRulesMinutes}{2.3}  % results: experiments/results/bench/2026-08-26T20-32-48Z_bench_6457ac7/benchmark.json
\newcommand{\vNgramMinutes}{2.3}  % results: experiments/results/bench/2026-08-26T20-32-48Z_bench_6457ac7/benchmark.json
\newcommand{\vLatencyRatio}{5037}  % results: experiments/results/bench/2026-08-26T20-32-48Z_bench_6457ac7/benchmark.json
\newcommand{\vFamilySize}{33}  % results: experiments/results/analysis/2026-08-26T21-20-11Z_p6-analysis_6457ac7/analysis.json
\newcommand{\vClusters}{10}  % results: experiments/results/analysis/2026-08-26T21-20-11Z_p6-analysis_6457ac7/analysis.json
\newcommand{\vArrangements}{1024}  % results: experiments/results/analysis/2026-08-26T21-20-11Z_p6-analysis_6457ac7/analysis.json
\newcommand{\vMinAttainableP}{0.001953}  % results: experiments/results/analysis/2026-08-26T21-20-11Z_p6-analysis_6457ac7/analysis.json
\newcommand{\vAlpha}{0.05}  % results: experiments/results/analysis/2026-08-26T21-20-11Z_p6-analysis_6457ac7/analysis.json
\newcommand{\vFalseDiscoveryRate}{0.05}  % results: experiments/results/analysis/2026-08-26T21-20-11Z_p6-analysis_6457ac7/analysis.json
\newcommand{\vPracticalThreshold}{0.02}  % results: experiments/results/analysis/2026-08-26T21-20-11Z_p6-analysis_6457ac7/analysis.json
\newcommand{\vImprovements}{11}  % results: experiments/results/analysis/2026-08-26T21-20-11Z_p6-analysis_6457ac7/analysis.json
\newcommand{\vRegressions}{5}  % results: experiments/results/analysis/2026-08-26T21-20-11Z_p6-analysis_6457ac7/analysis.json
\newcommand{\vBelowThreshold}{0}  % results: experiments/results/analysis/2026-08-26T21-20-11Z_p6-analysis_6457ac7/analysis.json
\newcommand{\vNoChange}{17}  % results: experiments/results/analysis/2026-08-26T21-20-11Z_p6-analysis_6457ac7/analysis.json
\newcommand{\vUnderpowered}{0}  % results: experiments/results/analysis/2026-08-26T21-20-11Z_p6-analysis_6457ac7/analysis.json
\newcommand{\vCrossCheckChecked}{33}  % results: experiments/results/analysis/2026-08-26T21-20-11Z_p6-analysis_6457ac7/analysis.json
\newcommand{\vCrossCheckDisagreements}{0}  % results: experiments/results/analysis/2026-08-26T21-20-11Z_p6-analysis_6457ac7/analysis.json
\newcommand{\vOldClusters}{5}  % results: experiments/results/analysis/2026-08-26T02-53-59Z_analysis_336175f/analysis.json
\newcommand{\vOldArrangements}{32}  % results: experiments/results/analysis/2026-08-26T02-53-59Z_analysis_336175f/analysis.json
\newcommand{\vOldMinAttainableP}{0.0625}  % results: experiments/results/analysis/2026-08-26T02-53-59Z_analysis_336175f/analysis.json
\newcommand{\vOldFamilySize}{11}  % results: experiments/results/analysis/2026-08-26T02-53-59Z_analysis_336175f/analysis.json
\newcommand{\vOldUnderpowered}{11}  % results: experiments/results/analysis/2026-08-26T02-53-59Z_analysis_336175f/analysis.json
\newcommand{\vCorpusForms}{960}  % results: corpus/manifest.json
\newcommand{\vCorpusFields}{10351}  % results: corpus/manifest.json
\newcommand{\vCorpusTemplates}{40}  % results: corpus/manifest.json
\newcommand{\vCorpusLocales}{6}  % results: corpus/manifest.json
\newcommand{\vCorpusFamilies}{5}  % results: corpus/manifest.json
\newcommand{\vCorpusTiers}{4}  % results: corpus/manifest.json
\newcommand{\vCorpusSeed}{20260825}  % results: corpus/manifest.json
\newcommand{\vTrainForms}{500}  % results: models/dev_metrics.json
\newcommand{\vTrainRows}{5402}  % results: models/dev_metrics.json
\newcommand{\vDevForms}{120}  % results: models/dev_metrics.json
\newcommand{\vDevRows}{1596}  % results: models/dev_metrics.json
\newcommand{\vExcludedForms}{100}  % results: models/dev_metrics.json
\newcommand{\vClassesFitted}{42}  % results: models/dev_metrics.json
\newcommand{\vDevMacroF}{0.7866}  % results: models/dev_metrics.json
\newcommand{\vDevAccuracy}{0.8421}  % results: models/dev_metrics.json
\newcommand{\vEceBefore}{0.0281}  % results: models/dev_metrics.json
\newcommand{\vEceAfter}{0.0454}  % results: models/dev_metrics.json
\newcommand{\vMacroFBeforeCalibration}{0.7690}  % results: models/dev_metrics.json
\newcommand{\vSwitchoverCount}{100}  % results: models/dev_metrics.json
\newcommand{\vParityRows}{1596}  % results: models/dev_metrics.json
\newcommand{\vParityAgreements}{1596}  % results: models/dev_metrics.json
\newcommand{\vParityDelta}{0.000000477}  % results: models/dev_metrics.json
\newcommand{\vFeatureWidth}{7986}  % results: models/train_manifest.json
\newcommand{\vChosenC}{8.0}  % results: models/train_manifest.json
\newcommand{\vOnnxOpset}{26}  % results: models/train_manifest.json
\newcommand{\vTauHigh}{0.8964}  % results: src/autofill_audit/audit/thresholds.json
\newcommand{\vTauLow}{0.6278}  % results: src/autofill_audit/audit/thresholds.json
\newcommand{\vTargetPrecision}{0.98}  % results: src/autofill_audit/audit/thresholds.json
\newcommand{\vDevPrecisionAtTau}{0.9954}  % results: src/autofill_audit/audit/thresholds.json
\newcommand{\vDevRecallAtTau}{0.3701}  % results: src/autofill_audit/audit/thresholds.json
\newcommand{\vRuleTauHigh}{0.70}  % results: src/autofill_audit/audit/thresholds.json
\newcommand{\vRuleTauLow}{0.40}  % results: src/autofill_audit/audit/thresholds.json
\newcommand{\vLlmTauHigh}{0.0}  % results: src/autofill_audit/audit/thresholds.json
\newcommand{\vLlmTauLow}{0.0}  % results: src/autofill_audit/audit/thresholds.json


% One defect, in the four parts that make a defect report useful to somebody
% who did not hit it: what was seen, how it was diagnosed, what was done, and
% what else was considered and refused. The fourth part is the one that is
% usually missing and is usually the reason the entry was worth writing.
\newenvironment{defect}[1]{%
  \par\medskip\noindent\textbf{#1}\par\nopagebreak\smallskip%
  \begin{list}{}{%
    \setlength{\leftmargin}{1.2em}%
    \setlength{\labelwidth}{0pt}%
    \setlength{\labelsep}{0pt}%
    \setlength{\itemindent}{0pt}%
    \setlength{\listparindent}{0pt}%
    \setlength{\parsep}{2pt}%
    \setlength{\itemsep}{2pt}%
  }\item[]%
}{%
  \end{list}\par\medskip%
}

\newcommand{\symptom}[1]{\par\noindent\emph{Symptom.} #1}
\newcommand{\diagnosis}[1]{\par\noindent\emph{Diagnosis.} #1}
\newcommand{\resolution}[1]{\par\noindent\emph{Fix.} #1}
\newcommand{\instead}[1]{\par\noindent\emph{Considered and refused.} #1}

\title{%
  \Large\textbf{autofill-audit: the debug report}\\[.6em]
  \large What went wrong, how it was found, and what was done about it\\[.4em]
  \normalsize Grouped by theme rather than by date
}
\author{Olajide Badejo}
\date{2026-08-27}

\begin{document}

\maketitle

\begin{abstract}
\noindent
This is the record of the defects, surprises and reversed decisions of the
project, drawn from an engineering log that was written as the work happened
rather than reconstructed afterwards. It is grouped by theme, because
chronology is how the problems were met and theme is how they are useful to
anyone else.

Every entry states the symptom, how it was diagnosed, what was done, and what
else was considered and refused. The fourth of those is the part that is
ordinarily missing from a defect record and is usually the reason the entry was
worth keeping: a fix without its rejected alternatives is a fact, and a fix with
them is an argument somebody can disagree with.

The document is a first class deliverable rather than an appendix. The
honest record of what went wrong is the part of a project that ordinarily
evaporates, and almost nobody publishes one.
\end{abstract}

\tableofcontents

\chapter{DOM traversal and shadow boundaries}

The traversal itself was right almost immediately. Everything in this chapter is
a case where the DOM said something true and the code drew the wrong conclusion
from it.

\begin{defect}{A label attached itself to the wrong control}
\symptom{On a fixture with two elements carrying the same identifier, a control
with no label of its own was reported with somebody else's label text, and the
classification that followed was confident and wrong.}
\diagnosis{The first implementation built a map from identifier to label and
handed the same label to every element carrying a duplicated identifier. That is
not what the platform does. A label's \ident{for} attribute names exactly one
element, the one the document's own lookup would return, and a duplicated
identifier does not make two elements labelled. It makes one labelled and one
not.}
\resolution{Resolve the association the way the platform resolves it: first
element wins, everything else has no label from that source and falls through to
the next source in the precedence order.}
\instead{Reporting a duplicate identifier as a defect and refusing to associate
either. That is a real finding and it is emitted, but it is a warning about the
page rather than a reason to withhold a label from the element that legitimately
owns it.}
\end{defect}

\begin{defect}{The canvas rule would have unbalanced every hostile checkout}
\symptom{The specification says a canvas rendered form like widget is reported
as undetectable rather than guessed at. Implemented literally, that emitted one
synthetic descriptor per large canvas, and every hostile checkout form in the
corpus then reported one more field than its answer key held. The extraction
reconciliation the phase gate asks for could never balance.}
\diagnosis{The specification describes a page whose form \emph{is} a canvas. The
corpus contains pages with plenty of real controls that also happen to carry a
canvas, which is a different shape, and the rule was being applied to both.}
\resolution{A canvas is always recorded as a region, on every page. The page
level undetectable descriptor is emitted only when the page yielded no controls
at all. The information is present either way; what changes is whether it counts
as a field.}
\instead{Excluding synthetic descriptors from the reconciliation. That would have
made the gate pass by making the gate weaker, and the gate exists to catch
exactly the class of error where the extractor and the answer key disagree about
how many fields a page has.}
\end{defect}

\begin{defect}{The frame accessibility test was asking the wrong question}
\symptom{The plan assumed the browser automation layer would fail to evaluate
inside a cross origin frame, and that this failure was how the tool would detect
one.}
\diagnosis{It does not fail. The automation layer drives each frame through its
own session, so it can evaluate inside a cross origin frame perfectly well. The
thing that makes such a frame unreadable to \emph{this tool} is not the
automation layer. It is the page's own same origin policy, and a tool that reads
across it would be reporting on a document the page cannot see either.}
\resolution{Make the accessibility test from inside the parent document, by
asking for the frame's document in a guarded block. That is exactly the check a
script on that page would make, and it is what the specification means by
\emph{cannot be accessed}.}
\instead{Using the automation layer's cross frame access and reporting on the
hosted document. It would have produced more findings and every one of them
would have been about a document the page's own author cannot reach from their
markup, which is advice nobody can act on.}
\end{defect}

\begin{defect}{Two offline fixtures for a problem that seemed to need a server}
\symptom{Testing frame behaviour appeared to require two origins, which appeared
to require a network or at least two local servers, which the testing rules
forbid.}
\diagnosis{Origins are cheaper than they look. A frame declared with inline
document markup inherits its parent's origin and so is same origin even under a
file scheme, where two separate local files are opaque origins to each other and
would not be. A frame loaded from a data address gets an opaque origin and so is
unreadable.}
\resolution{Two fixtures, both offline, both reproducing their case exactly, with
the reasoning recorded in the environment document so nobody replaces them with
a server later.}
\instead{A local server fixture in the test suite. It would have worked and it
would have made the whole suite depend on a port being free, which is the kind
of dependency that turns a green suite into an intermittently red one on somebody
else's machine.}
\end{defect}

\begin{defect}{One selector, four times, for a group of four radios}
\symptom{A radio group produced the same selector once per member.}
\diagnosis{The selector preference order includes a rank that selects on the
control's name attribute. Every member of a radio group shares one name, so that
rank produces one selector for the whole group and hands it to each member. The
sketch of the preference order in the specification does not mention this,
because the case does not arise for the controls it was written against.}
\resolution{Rank two requires the name to be unique within the form. A group
falls through to the next rank, which distinguishes its members.}
\instead{Special casing radio groups in the selector generator. The uniqueness
requirement is the general statement of the same rule and it covers checkbox
groups and any other shared name case without a second branch.}
\end{defect}

\begin{defect}{A child combinator that matched nothing}
\symptom{The second rank of the selector order, written literally from the
specification's sketch as a form element followed by a child combinator,
resolved to no element on every form the generator emits.}
\diagnosis{A child combinator matches only a control that is an immediate child
of the form element. Every form this generator emits, and essentially every form
on the web, wraps its controls in layout containers.}
\resolution{A descendant combinator. The deviation is recorded in the selector
module and in the handoff notes, because the extractor has to reproduce the
generator's selector convention exactly rather than approximately, and a
difference between the two would silently lose fields.}
\instead{Emitting forms without layout containers so that the literal reading
works. That would make the corpus unrepresentative of the thing being measured
in order to preserve a sketch.}
\end{defect}

\begin{defect}{Field order disagrees on exactly one tier, and it is not a defect}
\symptom{Reconciling extracted descriptors against answer keys showed an order
difference on the mixed markup tier and nowhere else.}
\diagnosis{The extractor returns document order, which the traversal
specification requires. An answer key lists its fields in the generator's slot
order, which appends a tier's extra controls to the end of the list. On the mixed
tier a hostile block sits in the middle of a form whose extra controls belong to
that block, so they render in the middle and the key lists them last.}
\resolution{The reconciliation compares selectors as a set and reports order
differences separately. The selectors are the contract; the order of the list is
not. That every affected form is on the mixed tier is asserted rather than
assumed, so a future order difference somewhere else fails the check.}
\instead{Sorting both sides before comparing. That would have hidden a real
future defect: an extractor that returned fields in an order unrelated to the
document would pass a sorted comparison forever.}
\end{defect}

\begin{defect}{Five signals with nowhere to go}
\symptom{The signal collection specification names several signals that the
descriptor schema has no field for.}
\diagnosis{Two sections of the specification disagree, and both are load bearing:
the collection list is what the traversal should gather, and the descriptor is
the fixed downstream contract that every golden test inherits.}
\resolution{Collect them, keep them on the raw record, and stop before the
descriptor. One of them, whether a select accepts multiple values, turned out to
be genuinely useful: a multi select with twelve options is a month picker to a
careless detector and a multiple choice control to a careful one.}
\instead{Widening the descriptor. That is a schema change that invalidates every
committed fixture expectation, and it should be made when something needs the
field rather than when something collects it.}
\end{defect}

\chapter{Waiting, flake, and process isolation}

The most expensive twenty minutes of the project were spent looking in the wrong
place because an error message named a symptom that had nothing to do with its
cause. That entry is first.

\begin{defect}{An error message about the wrong API}
\symptom{Adding a session scoped browser fixture to the audit test suite made the
extractor test suite fail with a message saying the synchronous browser API was
being used inside an asynchronous event loop, and recommending the asynchronous
API instead. Nothing anywhere was using the asynchronous API.}
\diagnosis{The synchronous wrapper drives a coroutine scheduler on one event loop
per thread. A second live browser session in the same thread finds that loop
already running and reports it in the only vocabulary it has, which is the
vocabulary of the asynchronous API. The message is accurate about the state it
found and useless about how the state arose.}
\resolution{One browser fixture in the root test configuration, shared by every
suite.}
\instead{Making each suite's fixture function scoped. That would have paid a
browser launch per test, and it would have left the same landmine for the next
person who added a second session concurrently rather than sequentially.}
\end{defect}

\begin{defect}{The end to end tests could not run the tool}
\symptom{Following on from the entry above: the command line tool opens a browser
of its own, and it cannot do that while the test session's browser is alive in
the same thread.}
\diagnosis{Same root cause, different consequence. It is not a defect in the
tool; it is a property of the synchronous wrapper.}
\resolution{The end to end tests run the tool in a subprocess. That turned out to
be an improvement rather than a workaround: a subprocess exercises the real exit
codes as a shell sees them, which is the actual contract the tool publishes, and
an in process call would have tested a function's return value instead.}
\instead{Running the end to end tests in a separate thread. It would have worked
and it would have tested a code path no user ever takes.}
\end{defect}

\begin{defect}{Choosing waiting constants that are neither flaky nor slow}
\symptom{Not a defect so much as a set of decisions that would have become
defects if they had been guessed. A page that never settles must not hang the
tool, and a page that is merely slow must not be cut off mid load.}
\diagnosis{Each constant is answering a different question and they should not be
derived from one another.}
\resolution{Each is a named constant with the reasoning beside it:

\begin{itemize}
  \item A generous load timeout, because a slow page is a real page, and a bound
    on it, because a page that never loads must not hang a command line tool.
  \item A network idle wait that gives up rather than failing, because plenty of
    healthy pages hold a socket open forever and giving up on idle is not an
    error.
  \item A DOM quiet period longer than an animation frame and shorter than the
    gap a page leaves before injecting a field it means a user to see.
  \item A total settle budget, so that an animation which mutates forever cannot
    hold the quiet period off indefinitely.
  \item A cap on controls, well above any real form, reported when reached and
    never silent.
  \item A frame depth bound, because three deep is an ordinary advertising stack
    and eight is pathological.
\end{itemize}}
\instead{One timeout for everything, tuned until the suite went green. That is
how a project acquires a constant nobody can explain and nobody dares change.}
\end{defect}

\begin{defect}{A window that would expire}
\symptom{A card expiry select has to offer years near the present, and the
extractor detects an expiry pair against the current year. A hardcoded window
would have made the corpus regenerate differently once the year turned, and the
reachability check and the committed sample would both have gone red on a
calendar boundary with no code change.}
\diagnosis{The corpus is required to be a deterministic function of its seed. A
clock is not a function of a seed.}
\resolution{The base year is a generator argument defaulting to the current year
and recorded in the manifest. The corpus is a deterministic function of the seed
and the base year together, and passing the manifest's base year back reproduces
the bytes exactly, indefinitely. The reachability check and the committed sample
both pin it.}
\instead{Freezing the year to a literal. The corpus would then have been
reproducible and progressively less realistic, and the expiry detection it exists
to exercise would have stopped being exercised.}
\end{defect}

\begin{defect}{The whitespace hook ate the golden snapshots}
\symptom{The first commit of the terminal renderer's golden snapshots rewrote
every one of them and left the test suite failing against its own repository.}
\diagnosis{The terminal rendering library pads a table row out to the full
terminal width. The pre-commit trailing whitespace hook trims it. Both are
behaving correctly and the combination is a snapshot that can never match.}
\resolution{The hook skips the golden directory. That is a narrow exemption for
the one directory whose entire purpose is to be byte exact.}
\instead{Stripping trailing whitespace when writing and when comparing
snapshots. That hides a real class of rendering defect, which is a renderer that
starts emitting different padding, and the whole point of a golden snapshot is to
notice exactly that.}
\end{defect}

\chapter{Text normalisation and the rule vocabulary}

The unglamorous layer, and the one where the only genuinely subtle defects of the
extractor lived. Two of the entries below were found by a property test in about
two seconds each, and neither would have been found by reading the code.

\begin{defect}{The published step order destroys its own worked example}
\symptom{Normalisation is specified as compatibility normalisation, then case
folding, then camel case splitting, then delimiter splitting, then a stoplist.
Implemented in that order, the specification's own worked example becomes
unreachable: the input folds to a single lowercase run in which no case boundary
exists, so the boundary step has nothing to find.}
\diagnosis{Two steps are in the wrong order relative to each other. Everything
the fold is specified for is a per character property, so it does not need to
happen before the boundary pass; the boundary pass genuinely does need the case
information the fold destroys.}
\resolution{Compatibility normalisation, then boundary insertion, then fold each
resulting token. Nothing the fold exists for is lost. What is gained is that the
boundary step works at all.}
\instead{Following the order literally and adjusting the worked example. That
would have been transcription rather than reading, and the resulting tokeniser
would have been useless on the camel case identifiers that are most of what a
hostile page leaves standing.}
\end{defect}

\begin{defect}{Combining marks were being treated as separators}
\symptom{A property test asserting that normalisation is idempotent failed
almost immediately, on the Turkish dotted capital letter I.}
\diagnosis{The splitter was built on the language's word character class, which
excludes combining marks. Folding the Turkish letter yields a base letter
followed by a combining dot; splitting on that dot drops the dot; a second pass
over the same text therefore returns something different from the first.}
\resolution{Treat combining marks as part of the token they attach to.}
\instead{Narrowing the property to the six locales the corpus covers and calling
idempotence a practical guarantee. This is the entry that most nearly went the
wrong way, and the reason it did not is worth keeping. Splitting on combining
marks shreds Devanagari, Thai, pointed Hebrew and anything in decomposed form.
None of that appears in this project's corpus, so the narrowed property would
have passed forever, and the defect would have been waiting for the first person
to point the tool at a page outside those six locales. The strategy went the
other way instead and now sweeps the whole basic multilingual plane.}
\end{defect}

\begin{defect}{Case folding does not always produce lowercase}
\symptom{The same property test failed again on the next run, on a different
input.}
\diagnosis{Cherokee folds upward. Folding a lowercase Latin letter beside a
Cherokee capital therefore leaves a lowercase letter next to an uppercase one:
a case transition that the camel case pass, which runs before the fold, never
saw. Feed the output back in and it splits into two tokens, which is a different
answer from the first.}
\resolution{Bracket the fold with boundary insertion rather than running it once
before. Folding before the boundary pass is not an option, because that is the
defect in the entry above, so the only order satisfying both constraints is to
insert boundaries, fold, and insert boundaries again.}
\instead{Excluding scripts that fold upward. There is one obvious such script and
assuming it is the only one is exactly the assumption that produced this defect
in the first place.}
\end{defect}

\begin{defect}{A stoplist entry that could never match}
\symptom{The one entry on the stoplist that names a real framework's generated
identifier prefix never fired.}
\diagnosis{It cannot survive the letter to digit split, which turns it into a
letter fragment and a digit fragment. A stoplist applied only after splitting
would never see it in the form the list is written in.}
\resolution{Apply the stoplist twice, to the delimiter separated fragments first
and to the finished tokens second, from one list.}
\instead{Adding the split halves to the list. The halves are both extremely
common substrings and adding them would have removed real signal from unrelated
identifiers.}
\end{defect}

\begin{defect}{A German placeholder read as a whole name field}
\symptom{On the hostile tier the label is stripped and the placeholder is all
that survives. The German email placeholder is the conventional documentation
address, the at sign does not survive normalisation, so the token stream reads as
three ordinary words, and the whole name rule fired on the first of them.
Ten confident wrong criticals across the hostile slice.}
\diagnosis{The rule table had no way to recognise a placeholder that is an email
address once the punctuation had been removed.}
\resolution{A rule matching the reserved documentation domain of the relevant
standards document, which is what essentially every placeholder on the web uses.}
\instead{Removing the whole name rule's weakest pattern, or special casing the
German locale. The domain rule turned ten wrong criticals into ten right ones
across the entire hostile slice rather than fixing the German case, because the
same placeholder convention is used in every locale in the corpus.}
\end{defect}

\begin{defect}{A Japanese section heading read as a card verification code}
\symptom{A login password field on a Japanese form picked up a card security code
classification from its surrounding context.}
\diagnosis{The rule for the Japanese security code matched the bare word for
\emph{security}, which is what a form's own security section heading contains.
The context tier then fed the heading to the rule.}
\resolution{Narrow the rule to the full compound.}
\instead{Removing the context tier for that label. The context tier is what makes
the hostile tier tractable at all, and removing it to fix one over match would
have cost far more than it saved. This defect and the one above were both found
by sweeping the generated corpus and cross checking every declaration finding
against the generator's own provenance, which is a check worth keeping.}
\end{defect}

\begin{defect}{A rule that cannot fire, kept anyway}
\symptom{The specification names the Japanese postal mark among the strings the
postal code label should match. It never fires.}
\diagnosis{Normalisation treats a lone symbol character as a delimiter, so the
mark never reaches a token stream. The Japanese postal field is reached through
its word form instead, which is why nothing was ever visibly broken.}
\resolution{The rule stays in the table, with a comment saying it cannot fire and
why, and a test asserting both halves of that: that the mark does not survive
normalisation, and that the field is still reached by the word form.}
\instead{Deleting the rule, or making it reachable. Deleting it would lose the
record that the specification asks for something this normalisation cannot
deliver. Making it reachable means changing what counts as a token character,
which changes every golden snapshot and every committed fixture expectation, and
that is a deliberate change with a cost rather than a quiet fix. It is written up
as future work rather than done under cover of a different phase.}
\end{defect}

\begin{defect}{The tie is the feature}
\symptom{Not a defect. It is recorded here because the first design was heading
for one and the correction changed the shape of the whole rule table.}
\diagnosis{German writes one phrase for both a whole street address and the first
of several address lines. French does the same. A field labelled only
\emph{Password} is a login field on one page and a registration field on the
next. Writing one pattern per label and hoping they do not overlap produces, in
each of those cases, a confident answer to a question the evidence does not
settle.}
\resolution{Make the overlaps deliberate. Both labels carry the same phrase at
the same weight, the tier ties by construction, and the engine falls through to a
tier that might separate them. A tie surviving every tier is
\ident{UNKNOWN}. The property that a correct page produces no finding above the
lowest severity fell out of this rather than being tuned into existence, and it
passed the first time it was run across the entire correct markup slice.}
\instead{Breaking each tie by picking a favourite. It would have looked like
better coverage on every metric except the one that matters, and it would have
produced a confident instruction to write the wrong token into somebody's
production markup.}
\end{defect}

\chapter{Generator determinism}

Same seed, same bytes. The gate is the generator run twice with a recursive diff
of the two outputs, and it is a gate rather than a test because a determinism
defect that reaches the corpus invalidates every measurement taken afterwards.

Two of the entries below are defects that had not yet caused a wrong byte and
would have, and those are the dangerous kind.

\begin{defect}{A shared value provider was a latent determinism defect}
\symptom{None yet, which is the point.}
\diagnosis{The first implementation cached one value provider per locale and
reseeded it per form. That is safe only while exactly one provider is alive at a
time. Nothing in the call path broke that invariant, which is what made it
dangerous: it would have broken the same seed same bytes rule silently, the first
time a caller built two forms concurrently, and the failure would have appeared
as a corpus that did not reproduce rather than as an exception anybody could
trace.}
\resolution{Each provider owns its own instance. Constructing one costs well
under a millisecond, so the whole grid pays a fraction of a second for an
invariant that cannot be broken from outside.}
\instead{Documenting the constraint and relying on callers to respect it. A
documented invariant that only one code path currently satisfies is an invariant
waiting for a second code path. Found by a test asserting that two providers
built with the same seed produce the same value, which they did not.}
\end{defect}

\begin{defect}{The mixed tier could produce a form with no clean section}
\symptom{Two templates out of the twenty five then in the grid produced
uniformly hostile forms carrying a mixed tier label.}
\diagnosis{The mixed tier draw forces at least one hostile section and at least
one clean one. The clean pass looked for a section that was not already hostile,
and when every section had drawn hostile there was no such candidate, so it
silently did nothing.}
\resolution{A section may now be promoted to clean unless it is the only hostile
one, and because a template is required to carry at least two sections there is
always a legal choice.}
\instead{Redrawing until a legal configuration appears. That changes the number
of draws taken, which changes the state of the generator, which changes every
subsequent form. A repair that consumes a variable amount of randomness is a
determinism defect wearing a fix's clothing.

The frequency is the part worth keeping: two templates in twenty five is exactly
the rate at which a defect gets shipped. The parametrised test covers all
templates rather than a sample, which is how it was caught.}
\end{defect}

\begin{defect}{Infinite recursion comparing two elements}
\symptom{A stack overflow while resolving a positional selector in the small HTML
tree used to check that generated selectors resolve.}
\diagnosis{Every element in that tree carries a pointer to its parent, and the
element type was a plain data class with a generated equality method. Resolving a
positional selector asks a list for the index of an element, which compares
elements, which compares parents, which compares children, and so on up and down
the tree.}
\resolution{Elements are identity objects. Equality generation is switched off
and the reason is in the class docstring, so nobody switches it back on for
convenience.}
\instead{Excluding the parent pointer from the comparison. That produces an
equality that says two structurally identical elements in different parts of the
tree are the same element, which is false and would have produced wrong selector
indices instead of a stack overflow. A crash is a better failure than a silently
wrong answer.}
\end{defect}

\begin{defect}{A module shadowed a data directory}
\symptom{Loading the locale tables looked one directory level too high and found
nothing.}
\diagnosis{The specification names the data directory with a plural noun, and the
obvious name for the module that loads it was the same plural noun with a source
extension, beside it. Importing the package qualified name then resolves to the
module rather than the directory, and asking the resource machinery for the
module's files walks to the package root.}
\resolution{The module is renamed, the directory keeps the name the specification
gives it, and the reason is in the module docstring so that nobody renames it
back on tidiness grounds.}
\instead{Renaming the directory. The directory name is specified and is part of
the layout other people read; the module name is not.}
\end{defect}

\begin{defect}{Two templates of one family drew the same values}
\symptom{Two structurally different templates of the same family, at the same
locale, tier and variant, produced identical sample values.}
\diagnosis{The specification writes the per form seed's input as the tuple of
family, locale, tier and variant. With several templates per family that tuple no
longer identifies a form.}
\resolution{The template identifier joins the digest input and the family stays
in it, so the input is a superset of the specification's rather than a
replacement for it.}
\instead{Replacing the family with the template identifier. Keeping both costs
nothing and means a future change to how templates are named cannot silently
collapse two families onto one seed.}
\end{defect}

\begin{defect}{Selector resolution checked without a browser}
\symptom{Not a defect. A design decision that prevents a whole class of them.}
\diagnosis{An answer key whose selectors do not resolve is worse than no answer
key at all: every downstream measurement silently loses the fields it could not
find, and the loss looks like a classifier failure rather than a corpus defect.
That check has to run in the phase that generates the keys, and that phase has no
browser.}
\resolution{A small resolver that handles exactly the subset of the selector
language this generator emits, and raises on anything else rather than quietly
returning no match for syntax it does not understand.}
\instead{Requiring a browser in the generator's gate. That would couple the
corpus phase's gate to the extractor phase's dependency and would stop the
reachability check from being able to check anything on a machine without a
browser installed.}
\end{defect}

\begin{defect}{A dependency classification that was reversed}
\symptom{The value provider library was recorded as a development dependency on
the reasoning that the generator runs from a checkout and never from an installed
package.}
\diagnosis{That was wrong, and it only became visible once the command surface
was wired up: corpus generation is a documented subcommand of the shipped
console entry point. A documented command that raises an import error on an
installed package is worse for a user than one extra dependency is for the
package.}
\resolution{The library moves to the runtime dependency set. Regenerating the
lock file afterwards produced a byte identical file, because the lock is compiled
with the development extra already and the resolver annotates both cases the same
way, so there is no lock commit. That is the right outcome rather than a skipped
step: regenerating a lock that does not change should not create a commit.}
\instead{A lazy import with a graceful refusal. That buys a smaller package at
the cost of a code path that only ever runs on a misconfigured install, which is
a code path nobody will test.}
\end{defect}

\chapter{Export fidelity and the featuriser fallback}

The phase that expected a fight with a converter and got a different problem
instead.

\begin{defect}{The fallback was the design, not a retreat}
\symptom{The plan anticipated a fight with the converter for text vectorizers and
named a fallback: move the whole text pipeline into hand written code, export
only the linear layer, and take that quickly rather than spending a day coaxing a
converter.}
\diagnosis{The fight never started, because the feature set was never
convertible. Three of the five feature blocks are categorical one hots, option
shape booleans and structural buckets, and no standard transformer produces
those. Putting the whole featurisation in the graph would have meant writing
custom converters for three bespoke transformers \emph{and} trusting the
vectorizer converters, which is a larger version of exactly the risk the fallback
exists to avoid.}
\resolution{The hand written featuriser, shared by training and inference. It is
the only design in which one implementation serves both paths, which the feature
specification required from the beginning.}
\instead{Writing the custom converters. The second argument against arrived for
free and will matter longer: a vectorizer anywhere in the inference path puts a
full machine learning framework on the dependency list of a tool that installs
with one command. The route taken means an installed package needs an array
library and a runtime and nothing else.}
\end{defect}

\begin{defect}{A reimplementation that is a claim rather than an identity}
\symptom{A hand written character n-gram analyser is slower than a compiled one,
and its faithfulness to the reference analyser is an assertion.}
\diagnosis{So assert it, in a test, rather than in a comment.}
\resolution{A test asserts, over inputs including empty strings, words shorter
than the n-gram length, and Japanese text, that this project's enumeration is
character for character the reference implementation's.}
\instead{Spot checking a few strings. The detail an independent reimplementation
gets wrong is the short word rule: a word shorter than the n-gram length yields
its padded self once rather than once per length. Getting that wrong silently
triples the weight of every two letter token, and it would not show up in a spot
check of ordinary words.}
\end{defect}

\begin{defect}{An opset probe, and a warning that is recorded rather than
suppressed}
\symptom{The export is specified to use the newest opset the resolved runtime
accepts, determined by probe. The probe walks down from the newest the installed
format library defines and takes the first that both converts and opens in a
session. The converter warns at every attempt that the requested opset is above
the newest it has been tested against.}
\diagnosis{The warning is not noise. It is precisely the case where a successful
export proves nothing: the converter is saying it has no test coverage for the
combination it just produced.}
\resolution{The warning is recorded in the decision record rather than filtered,
and the thing that proves something is the parity gate. That is why the export
phase has a gate rather than a deliverable.}
\instead{Suppressing the warning and trusting the successful conversion. A
conversion that succeeds and produces different numbers is the failure mode the
whole gate exists for.}
\end{defect}

\begin{defect}{The graph cannot read its own weights}
\symptom{Naming the n-grams behind a prediction, which is what the evidence
requirement in this project's first law demands, could not be done from the
inference path.}
\diagnosis{The exported graph is one classifier node followed by a normaliser,
and the coefficients, the intercepts and the class order live in the node's
\emph{attributes} rather than in graph initialisers. The runtime does not hand a
caller a node's attributes. Reading them would have required the full model
format library at runtime to produce one line of a report.}
\resolution{A separate evidence table, holding each class's highest weighted
features, written by the same training run that produces the model.}
\instead{Adding the model format library as a runtime dependency, or dropping the
evidence line from findings. The first inflates every install for one line of
output; the second violates the law the tool is built around.

A derived file is a drift risk, so the parity suite reads the weights back out of
the graph and asserts that the table agrees with them. A stale evidence file now
fails the build rather than naming the wrong n-grams in a finding, which would
have satisfied the evidence law in form and violated it in substance.}
\end{defect}

\begin{defect}{The first training run was thrown away}
\symptom{The manifest of the first full training run recorded that the working
tree was dirty.}
\diagnosis{The reproducibility rules say a result produced from a dirty tree is
marked dirty and is not citable. The flag did exactly what it exists for.}
\resolution{The run was discarded and repeated from a clean tree.}
\instead{Citing it with a note. The flag needed one fix of its own to be honest:
the output directory is untracked on a first run, so counting it would have
marked every first training run dirty by construction and made the flag mean
nothing. It now excludes the output directory and only that.

This entry has a sequel. See the open item in
Chapter~\ref{ch:still-open} about the training manifest of the currently shipped
model, which is marked dirty for a different and less avoidable reason and was
not thrown away.}
\end{defect}

\chapter{Calibration and threshold derivation}

The phase where most of the work turned out to be arranging for numbers to be
checkable rather than producing them.

\begin{defect}{A threshold definition whose answer is always zero}
\symptom{The low threshold was specified as the smallest confidence at which the
near miss band still captures at least half of the recall the high threshold
gives up. Written down carefully, that has a degenerate answer.}
\diagnosis{Captured recall only ever falls as the threshold rises, so every value
below a qualifying one also qualifies, the smallest qualifying value is always
zero, and a threshold of zero says nothing.}
\resolution{The non degenerate reading is the \emph{greatest} qualifying value,
and it was written into the prediction file before the model existed.}
\instead{Picking whichever reading produced a nicer boundary once both were
computable. Writing the reading down first settled a real ambiguity at a point
where there were not yet two candidate answers to choose between, which is the
only time that choice is free.}
\end{defect}

\begin{defect}{One threshold document would have moved every rule engine finding}
\symptom{The derived n-gram boundary sits far above the rule engine's confident
band. A single pair of thresholds shared by both engines would have moved every
rule engine finding on every page, churned every golden snapshot, and done it as
a side effect of training a model.}
\diagnosis{A confidence scale is a property of an engine. A calibrated
probability and a rule tier band are different quantities that happen to be
written in the same units.}
\resolution{The threshold document carries a block per engine at a bumped schema
version, the loader takes an engine name, and the command line resolves the
engine first and then asks for that engine's block. The rule block is byte
identical to what the previous phase committed, and the golden diff after the
whole phase is the version string and nothing else, which is the evidence that it
worked.}
\instead{One global pair, or an inheritance rule for an engine with no block. An
engine with no block raises. That was designed to bite when the language model
arrived, and it did, which is how that engine's threshold decision came to be
made deliberately and in advance rather than by inheriting a boundary derived on
a different scale.}
\end{defect}

\begin{defect}{A precision of one over a denominator nobody could see}
\symptom{The derived block reported that the precision target was met.}
\diagnosis{A precision of one over fifteen accusations and a precision of one
over fifteen hundred are the same number and are not the same claim. A reader of
the file could not tell which they were looking at.}
\resolution{The derived block records the denominators beside the rates, and the
model card says the same thing in words.}
\instead{Reporting the rate alone and mentioning the denominator in prose
somewhere else. The file is what the runtime reads and what a reader opens
first.}
\end{defect}

\begin{defect}{Calibration made the summary statistic worse}
\symptom{Expected calibration error was lower before calibration than after it
(\vEceBefore{} against \vEceAfter{}).}
\diagnosis{The uncalibrated model on this corpus is already close to honest.
Refitting a one versus rest sigmoid per class on a development split of one
template per family, and then renormalising, moves classes that were already well
behaved.}
\resolution{The calibrator stays, and the model card says the error went the
wrong way and by how much. Two reasons: macro-F1 is higher with it
(\vMacroFBeforeCalibration{} before against \vDevMacroF{} after), and the
argument for calibrating at all is about the number a single finding quotes to a
developer rather than about an average gap across a split.}
\instead{Dropping the calibrator. Removing a component because its summary
statistic got worse, without re-registering a prediction first, is tuning on the
thing being measured. The honest move was to keep the design, publish the
number that went the wrong way, and say why the design still stands.}
\end{defect}

\begin{defect}{The sweep chose the edge of its own grid}
\symptom{The chosen inverse regularisation is the largest value the
pre-registered grid offered (Table in the main report's classifier chapter).}
\diagnosis{Either the grid is too narrow or the model wants less regularisation
than expected, and the data cannot distinguish those without a wider grid.}
\resolution{The grid was not widened. The consequence is recorded in the model
card instead: this model may be less regularised than a wider grid would have
chosen.}
\instead{Widening the grid and re-running. Widening a pre-registered grid because
the answer landed at its edge is how a sweep becomes a search for a number. A
later phase that revisits it re-registers first.}
\end{defect}

\begin{defect}{An abstention that was right by accident}
\symptom{A model emitting the unknown label with a high calibrated probability
reached the correct outcome, but through the wrong mechanism: the declaration
lookup returned nothing, rather than the confidence gating anything.}
\diagnosis{Right answer, wrong reason, and a wrong reason is a defect waiting for
the code around it to change.}
\resolution{An explicit branch. When the winning class is the unknown label the
reported confidence is zero, exactly as the rule engine reports it, and the
calibrated probability of the runner up stays on its own field where the
confusion analysis can still read it.}
\instead{Leaving it, on the grounds that the outcome was correct. The outcome was
correct for a reason that no test asserted and that any refactor of the
declaration lookup would have silently removed.}
\end{defect}

\begin{defect}{The model's first wrong accusations}
\symptom{After the corpus repair, the n-gram engine accused
\vNgramMissingAccusations{} fields of a missing declaration and was wrong about
\vNgramMissingWrong{} of them. That is the first time in this project that an
engine told a developer to add a token the answer key disagrees with.}
\diagnosis{Not a regression in the threshold policy. The target precision behind
the high threshold is the pre-registered one and the derivation script was run
unmodified. What changed is that the derived threshold now sits on a model whose
development precision at that point was itself below one, so a test split
precision below one is the design working rather than failing.}
\resolution{Recorded rather than rounded, in the engineering log, the model card
and the main report. A precision target is a target, not a guarantee.}
\instead{Raising the target until the wrong accusations disappeared. That is
tuning a pre-registered constant against the test split, which is the single
worst thing available in this project's rulebook.}
\end{defect}

\begin{defect}{A class that finally reached the isotonic switchover}
\symptom{Not a defect. A prediction that was correct for one corpus and stopped
being correct for the next.}
\diagnosis{The earlier phase predicted that no class would ever have enough
development positives for isotonic regression, and for its corpus that was true
of every class. After the repair, the label covering search boxes and consent
checkboxes crossed the switchover.}
\resolution{The per class method is recorded in the model card, so the change is
visible rather than implicit.}
\instead{Nothing to refuse here. It is in this chapter because a prediction that
holds under one corpus and fails under a larger one is exactly the kind of thing
that gets quietly forgotten, and the record of it is what makes the next such
prediction more careful.}
\end{defect}

\chapter{Language model schema behaviour and serving}

The chapter whose most interesting entry is a failure mode that did not occur.

\begin{defect}{A seven and a half gigabyte model on disk and unreachable}
\symptom{The serving runtime was already installed on the host operating system
with the model pulled. From inside the Linux subsystem the model blob was
visible on the filesystem and the server was unreachable over the network.}
\diagnosis{The host installation binds to loopback only, and loopback inside the
subsystem is not the host's loopback.}
\resolution{Install the runtime inside the subsystem and point its model
directory at the host's store. The existing blob is reused with no re-pull.}
\instead{Re-pulling the model inside the subsystem, or reconfiguring the host
server's bind address. The first costs the download again and doubles the disk;
the second changes a setting on the developer's own machine outside the
repository, which is the kind of environment drift that makes a result
irreproducible for the next person.}
\end{defect}

\begin{defect}{A keep alive long enough to be convenient and short enough to give
the card back}
\symptom{A long keep alive holds three quarters of the accelerator for as long as
it lasts. A short one puts a model load into the next call.}
\diagnosis{These are the two failure modes and they are not symmetric: an
occupied accelerator is a nuisance, and a cold load inside a measured run is a
corrupted latency distribution.}
\resolution{A bounded keep alive of a few minutes, an explicit unload when a run
finishes, and every latency figure taken against a warmed session so that the
reload costs nothing methodologically as long as the warm up precedes the
timing.}
\instead{An unbounded keep alive. It would have made the measurement identical
and left the accelerator occupied indefinitely, which is a decision that belongs
to whoever owns the machine rather than to the benchmark.}
\end{defect}

\begin{defect}{The engine order in the benchmark is a methodological choice}
\symptom{The benchmark runs the language model first, which is not the order the
tables present.}
\diagnosis{The two classical engines take a couple of minutes each of browser
time with no inference call in them. Running the language model last would have
meant several idle minutes before its first call, against a keep alive of a few
minutes, so the keep alive could expire inside the measured run and put a cold
model load into the per field latency distribution.}
\resolution{Warm immediately before, run it first, and record the ordering and
its reason in the run's own notes so that the table's presentation order cannot
be mistaken for the execution order.}
\instead{Presenting the tables in execution order. The tables are ordered for a
reader; the execution is ordered for the measurement, and conflating the two
would have made one of them worse for no gain.}
\end{defect}

\begin{defect}{The obvious leak check is wrong in both directions}
\symptom{The declared attribute value must be dropped from what the model sees,
or it reads the answer off the page on every clean tier form. The obvious check
is a substring search for the declared token in the serialised payload.}
\diagnosis{That check has two failure modes and both are fatal. A field named
\ident{email} that also declares an email autofill token would trip it, and that
combination is common enough in a clean corpus that the check would have fired on
the first real run and been switched off within the hour. Meanwhile a declaration
reaching the payload through some derived value would not trip it at all.}
\resolution{Check \emph{independence} instead: pruning a descriptor and pruning
the same descriptor with its declaration stripped must produce identical
payloads. That has neither failure mode, and it is asserted on every request
rather than reviewed once.}
\instead{The substring check with an allow list of tokens that are also common
identifiers. The allow list grows until the check means nothing, which is the
same failure the traceability checker was designed to avoid.}
\end{defect}

\begin{defect}{The keep list is a floor, not a ceiling}
\symptom{The specification's list of what pruning keeps omits several signals
that the featuriser feeds to the n-gram model.}
\diagnosis{Withholding a signal one engine has from the other tilts the
comparison exactly as far as handing over the whole document would, in the
opposite direction and with a better conscience.}
\resolution{The rule applied is that the language model sees what the featuriser
sees and nothing more. The three named drops stand, because none of them reaches
the featuriser either. The deviation from a literal reading of the list is
recorded in the prompt module's docstring with the reason.}
\instead{Following the list literally. It would have been defensible by citation
and indefensible as a measurement, and the specification's own reason for having
the list is fairness rather than the list itself.}
\end{defect}

\begin{defect}{A threshold block that had to be unflattering}
\symptom{The threshold loader raises for an engine with no block, by design. The
language model needed two numbers and the specification forbids its self reported
confidence from feeding the threshold policy.}
\diagnosis{The only pair satisfying both constraints is zero and zero.}
\resolution{Zero and zero, registered in the prediction file before the run, with
both consequences written down in advance: the engine gets no threshold
protection and accuses on every committed prediction, and the low confidence
finding becomes unreachable for it.}
\instead{Deriving a boundary on the self reported scale against a precision
target, the way the n-gram engine's was derived. It would probably have raised
the engine's finding level precision substantially. That it would have flattered
the engine is exactly why it had to be refused before the measurement rather than
after, and why the decision is in a committed file with a timestamp rather than
in a paragraph written afterwards.}
\end{defect}

\begin{defect}{The failure mode the retry ladder exists for did not occur}
\symptom{Across \vLlmRequests{} requests, \vLlmRetried{} needed a retry and
\vLlmFailed{} failed after one. The registered prediction was that between none
and one in twenty would need a retry.}
\diagnosis{Passing the response schema with the request lets the server constrain
decoding, which appears to remove the failure mode rather than reduce it.}
\resolution{Report the zero. The retry ladder is still there and is still tested
against a dozen recorded transcripts including deliberately malformed ones, so
its behaviour is covered without a network and without a key; on this run it
never fired.}
\instead{Reporting robustness in prose, or removing the untriggered ladder. A
benchmark whose most interesting failure mode did not occur is a result, and the
honest report of it is the zero rather than a paragraph about how robust the
implementation is. Removing the ladder would leave the next serving route, which
may not constrain decoding, with no handling at all.}
\end{defect}

\begin{defect}{The prompt fingerprint that fails on purpose}
\symptom{Editing the system prompt fails a test that holds a digest of the
assembled prompt and schema.}
\diagnosis{Working as intended. The prompt is part of what a result means, so a
silent prompt edit would invalidate every committed language model result without
anything going red.}
\resolution{The correct response is to bump the prompt version, update the
literal, and record the change in the changelog. Updating the literal alone is
the one thing the test exists to prevent.}
\instead{Comparing the prompt loosely, or leaving it untested. The whole reason a
result file records a prompt version is so that two results can be compared, and
a version that does not change when the prompt does is worse than no version at
all.}
\end{defect}

\begin{defect}{Cost is null, and null is not zero}
\symptom{The cost column is empty for every engine in every table.}
\diagnosis{The two classical engines have no provider and the language model ran
locally. Nothing was billed.}
\resolution{Record null. Zero is a measurement and the electricity was not free;
null is the absence of one. The price table document in this repository
deliberately contains no prices, because no vendor page was retrieved during this
work and a remembered price is a fabricated number.}
\instead{Writing zero, or computing a list price estimate from a remembered rate
card. The second is worse than the first: it produces a plausible number with a
currency symbol on it, which is exactly the shape of claim that survives review.}
\end{defect}

\begin{defect}{The evaluation runner pays for the language model twice}
\symptom{The run made twice as many requests as there are forms, and about half
its wall clock is a second inference pass.}
\diagnosis{The evaluation runner classifies every form twice: once directly,
which is where the run log's predicted label and per field latency come from, and
once through the audit engine, which is where the finding codes come from. For
the two deterministic engines the passes agree by construction and cost
microseconds. For a language model it doubles the inference.

There is a second consequence and it is the one that matters to a reader of the
finding level tables. A row's predicted label and its finding codes come from two
different draws of a model that is not guaranteed to be deterministic. Measured
on this run, using the deterministic rule engine as a control for the check's own
edge cases, the excess disagreement is on the order of two fields in seven
hundred.}
\resolution{Not fixed. Documented in the main report's limitations chapter, in
this report, and in the future work list, with the measured magnitude.}
\instead{Fixing it in the phase that found it. Fixing it changes what the
evaluation runner does for every engine, and doing that after seeing the results
is the regeneration this project's reproducibility rules forbid. The fix is to
pass the first pass's predictions into the audit rather than letting the audit
reclassify, which also turns the two numbers into one fact, and it belongs to a
phase that re-registers its predictions and re-runs everything.}
\end{defect}

\chapter{Friction at the package boundary}

The evaluation harness is consumed here as an ordinary released dependency: not
vendored, not submoduled, not copied file by file. That separation is a
deliberate engineering claim, and the friction it exposes is the point rather
than an inconvenience. Vendoring the code would have destroyed the evidence by
making the boundary unobservable.

Every entry below has an issue filed on that repository with a concrete API
proposal, and every workaround that stands here in the meantime is labelled in
every result file it produces, because a workaround that is invisible in the
output is indistinguishable from a capability.

\begin{defect}{The anticipated friction was aimed one layer too low}
\symptom{The ledger predicted, before the integration began, that categorical
outcome support in the permutation machinery would be the problem. It was not
needed at all.}
\diagnosis{The p value estimator never sees an outcome. It sees an observed
effect and a null distribution, both numbers. The categorical part of this
project's problem lives entirely in the statistic, and the statistic is the
caller's. What actually broke was one layer above, in the data model.}
\resolution{Record the miss in the ledger rather than quietly replacing the
prediction with what turned out to be true.}
\instead{Rewriting the anticipation to match the outcome. A ledger that only ever
records correct anticipations is a ledger nobody can learn anything from.}
\end{defect}

\begin{defect}{The ingestion layer refuses a run log that has no time axis}
\symptom{The harness's parser claims a run log of this project's schema and then
refuses it, reporting that the records carry no step field and naming the several
spellings it would accept.}
\diagnosis{There is no step to add. The file has one row per classified field and
no time axis at all. Inventing a step from the row index would let the harness's
window comparison treat a set of unrelated fields as a stationary process, which
is a stronger and false claim.}
\resolution{Filed with a concrete proposal. In the meantime the bridge reads its
own file format, which is a few lines, and nothing is reshaped at the last moment
to fit an API that does not quite fit.}
\instead{Emitting a synthetic step column. It would have made the parser accept
the file and would have made every subsequent analysis silently wrong.}
\end{defect}

\begin{defect}{No public entry point accepts a cluster assignment}
\symptom{The resampling unit is non negotiable in this project's design, and
neither of the harness's comparison modes is clustered. Neither is paired, and
both reduce a metric series to a final window mean.}
\diagnosis{Three separate assumptions, all of them consequences of the harness
having been extracted from a program that observed one training run over time.}
\resolution{The clustered null is built on this side and handed to the harness's
p value estimator, which is the one part of the interface that accepts a caller
supplied null. The estimator and the correction stay on the far side of the
boundary. The manifest of every analysis records that this happened, in its notes
and in a block that names what the harness was and was not used for.}
\instead{Degrading silently to unclustered resampling. That would have made every
comparison look more significant than it is, and it was the one outcome ruled out
before the integration began.}
\end{defect}

\begin{defect}{A verdict attached to the wrong comparison, and nothing raised}
\symptom{The first cross check run attached the verdict for one metric to a row
labelled with a different one. Every value in the row was plausible.}
\diagnosis{The harness's classification entry point ends by returning its
findings in severity order rather than input order, and the tag on a finding is
not a unique key when one metric is compared across several slices, which is
exactly this project's family. A positional zip therefore attaches every verdict
to the wrong comparison, and because every value is plausible, nothing raises.}
\resolution{The bridge rejoins on the identity of the result object each finding
carries. Caught by reading one line of output that should have said one metric
name and did not.}
\instead{Sorting both sides by tag before zipping. The tag is not unique, so that
produces a different wrong answer. Filed as its own issue with two proposed
fixes, because joining on object identity is not a contract worth depending on
and the fix belongs on the other side of the boundary.}
\end{defect}

\begin{defect}{A relative gate where an absolute one was pre-registered}
\symptom{The harness's practical effect threshold is a relative percentage. This
project pre-registered an absolute difference in the metric.}
\diagnosis{A relative gate turns one pre-registered rule into a different rule in
every slice, because the same relative percentage means something different
against a large baseline than against a small one.}
\resolution{The practical gate is applied here, and the harness's classification
runs alongside as a visible cross check so that the two verdict columns can be
compared. On this data the gap changed no verdict:
\vCrossCheckDisagreements{} of \vCrossCheckChecked{} comparisons disagreed.}
\instead{Adopting the relative gate because it happens to agree on this dataset.
A gap that does not bite on one dataset is still a gap, and the dataset it would
bite on is any slice with a small baseline, which is most of what a wider
benchmark would add.}
\end{defect}

\begin{defect}{A truthful mode string the harness cannot label}
\symptom{The comparison result's mode field is an unvalidated string, so the
bridge writes what actually ran. Two of that type's derived members look the
value up in a table holding the two training modes and raise on anything else.}
\diagnosis{The type models a closed set of modes in its derived members and an
open one in its field.}
\resolution{The bridge does not call those two members and serialises what it
needs itself.}
\instead{Writing one of the two known mode strings so that both members work.
That would make the derived members work and would be a false statement about
which test was run, printed in the result file. This one is deliberately not
filed on its own: it is a consequence of the same modelling assumption as the
clustering issue, and if a paired clustered mode lands there it will bring its
own label with it.}
\end{defect}

\begin{defect}{No typing marker, so every value crossing the boundary is
unchecked}
\symptom{Strict type checking refuses to look inside the package, so every value
that crosses the boundary arrives untyped.}
\diagnosis{The source is thoroughly annotated. The marker file is the whole fix
and it is missing.}
\resolution{A configuration override in this project's checker settings, filed on
the other side, with the cost recorded: the bridge is the one module in the
source tree whose external calls are unchecked, which is another reason for it to
be the only module that makes them, and a test walks every source file and fails
if a second importer appears.}
\instead{Writing type stubs here. Stubs for somebody else's package, maintained
in a consumer's repository, go stale silently and would have hidden exactly the
kind of API change the boundary exists to expose.}
\end{defect}

\begin{defect}{A statistics dependency that drags in a logging stack}
\symptom{Installing the harness pulls more than a dozen transitive packages,
including a training metrics logging stack, a dataframe library, a plotting
library and a web application toolkit, for a consumer that uses three functions.}
\diagnosis{The harness's ingestion layer reads event files and its reports draw
plots, and neither is separable at install time.}
\resolution{The dependency sits in the development extra rather than in the
runtime dependencies. The consequence is stated in the packaging file, in the
ledger and in the main report rather than left to be discovered: \textbf{the
shipped tool cannot compute its own significance tests}.}
\instead{Putting it in the runtime dependencies. An install of a command line
auditor that dragged a training metrics logging stack onto a developer's machine,
so that a command they will never run could compute a permutation p value, would
be the wrong trade. The narrower extra that would fix it properly is proposed in
the ledger.}
\end{defect}

\begin{defect}{The part furthest from its origin transferred without a scratch}
\symptom{Not a defect. The most interesting result of the whole integration.}
\diagnosis{The step up correction procedure needed nothing at all. It is a pure
function over a list of p values, it is checked in that repository against the
worked example from the original paper, and it did the right thing here on the
first call. The p value estimator fit because it takes the null as an argument.
The verdict vocabulary already contained the three categories this project needs
plus a fourth for a design that cannot reach its own level, which turned out to
be the category every comparison in the first analysis landed in.}
\resolution{Say so, in the ledger and in the main report.}
\instead{Reporting only the friction. The split runs cleanly through the middle
of the package: the parts built around the \emph{statistics} generalised, and the
parts built around the \emph{shape of a training run} did not, because that shape
was never a statistical assumption in the first place. It was the shape of the
one program the package was extracted from, and that is a sharper finding than
either half alone.}
\end{defect}

\chapter{The power repair}

The largest single defect in the project was not in any line of code. It was in
the shape of the corpus, it was found by arithmetic rather than by a failing
test, and it invalidated the ability of an entire phase to certify anything.

\begin{defect}{A design that could not reach its own significance level}
\symptom{Every one of the \vOldFamilySize{} comparisons in the first analysis
came back inconclusive, several of them sitting exactly at the design floor,
meaning the observed difference was more extreme than every other arrangement the
design permits and was still not significance.}
\diagnosis{Pure arithmetic, and it needed no test split measurement at all. The
resampling unit is the template. Whole templates are assigned to partitions, and
the realised grid put one template per family in the test partition, so the test
split held \vOldClusters{} templates. A paired sign flip permutation over that
many clusters has \vOldArrangements{} arrangements, so the smallest attainable
two sided p value is \vOldMinAttainableP{}, and the pre-registered false
discovery level is \vAlpha{}.

\textbf{More resamples do not help.} At that cluster count the test is already
exhaustive. Only more clusters help.}
\resolution{Widen the corpus. The test partition now holds \vTestTemplates{}
templates, giving \vArrangements{} arrangements and a smallest attainable p of
\vMinAttainableP{}, and \vUnderpowered{} comparisons in the headline analysis are
reported as inconclusive for design reasons.}
\instead{Two other options were on the table and both were refused in writing.

\emph{Accept the floor and report every comparison as inconclusive.} That leaves
the headline experiment answering its question with effect sizes and no
significance at all, which is a legitimate paper and a poor engineering result.

\emph{Pre-register a different alpha.} A level one order of magnitude looser
would clear the floor. It would also have been chosen \emph{after} discovering
that the original level was unreachable, which is indistinguishable from moving
the goalposts however carefully it is worded.}
\end{defect}

\begin{defect}{Computing the power before the runs rather than after them}
\symptom{Not a defect. The decision that made the entry above reportable.}
\diagnosis{The arithmetic above was computed and written into the statistical
policy file \emph{before} the test split runs, from the split assignment alone.}
\resolution{Publish it as a finding.}
\instead{Discovering it after the runs came back inconclusive and writing it up
then. The two documents would have contained the same sentences and would have
had entirely different evidential value. Writing it afterwards would have been
indistinguishable from an excuse.}
\end{defect}

\begin{defect}{Every recorded number in the repository moved at once}
\symptom{Widening the corpus changes the manifest digest, which invalidates the
descriptor cache, which changes the training partition, which changes the model,
its vocabulary, its calibration and its evidence table, which changes both
derived thresholds, which changes every metric in the model card, and separately
changes the test partition, which moves both engines' test split numbers in
either direction.}
\diagnosis{All of it follows mechanically from one deliberate change, and none of
it is a regression. But a reader comparing the project against its own history a
month later would otherwise see every number change at once with no stated
cause.}
\resolution{The explanation, the changelog entry and the prediction file are one
commit and it is the \emph{first} commit of the phase, before the templates, the
regeneration and the retraining. Then the model artefacts and the model card in a
single commit, as the model card rule requires. Then the new runs and the new
analysis.}
\instead{Doing the work and writing the explanation at the end. The ordering of
the commits is the part that matters, because it is the part a reader can check.}
\end{defect}

\begin{defect}{Three labels the model had never seen}
\symptom{Three labels sat at zero rows in the training partition, which means the
model had been asked about classes it could only ever get wrong.}
\diagnosis{Nothing in the growth rule required a training floor. Every clause was
about existence somewhere in the system, not about presence in the partition the
model actually learns from.}
\resolution{A new clause: every label a model can predict must appear in the
training partition with at least a full template's worth of rows. The constant is
arithmetic rather than taste. A template is generated in six locales and four
tiers at one variant, and the held out locale is lifted out of training, so one
training template contributes exactly that many training rows for a field it
carries in every locale. The constant therefore says \emph{at least one whole
training template carries this label}, which is the smallest statement about a
label that is not an accident of one locale profile. The new templates were
designed to carry the starved labels deliberately rather than to acquire them by
luck.}
\instead{Dropping the three labels from the model's output space. They are
specification tokens that real forms use, and a tool that cannot name a field
because its own corpus was thin is a tool with a gap in the middle of its
taxonomy.}
\end{defect}

\begin{defect}{What deliberately did not move}
\symptom{Not a defect. The list is here because it is the part that makes the
repair honest.}
\diagnosis{A phase that moved a threshold and a corpus at the same time would
have produced two changes and no attribution.}
\resolution{The seed stays where it was set. The held out locale stays the same
locale. The excluded partition keeps its rule. The alpha, the false discovery
rate, the absolute practical effect threshold, the clustering unit and the target
precision behind the high threshold are all exactly what the earlier phases
pre-registered.}
\instead{Taking the opportunity to adjust anything else while the numbers were
going to move anyway. It is the most tempting moment in a project to do that, and
it is the moment at which a repair becomes a rewrite.}
\end{defect}

\begin{defect}{The old result files were kept}
\symptom{Two test split runs in the repository were taken on a corpus that no
longer exists in that shape, and their numbers are not comparable to anything
current.}
\diagnosis{They were correctly taken and correctly reported.}
\resolution{They stay, unedited. Result files are append only history. The honest
description of them is the first, underpowered measurement, and they are cited in
exactly one place for exactly one purpose, which is to show what the repair did
to each engine, with the caveat attached that the two columns are not two
measurements of the same thing.}
\instead{Deleting them, or regenerating them against the new corpus. The
reproducibility rules forbid regenerating to fix a number, and deleting them
would remove the only evidence that the project found its own design defect
rather than being told about it.}
\end{defect}

\chapter{The checks, and the checks on the checks}

Enforcement infrastructure has a failure mode of its own: a check that cannot
pass its own rule gets an exemption, and an exemption is where enforcement starts
to rot. Three entries below are about checks that had to be built so that they
could survive themselves.

\begin{defect}{The dash checker was its own first violation}
\symptom{The first draft of the checker that forbids two specific characters
declared those two characters as string literals, which made the file an
immediate violation of the rule it enforces.}
\diagnosis{A checker that names what it forbids will contain what it forbids.}
\resolution{Build the characters from their codepoints at import.}
\instead{Exempting the checker from its own scan. The same reasoning applies to
the commit message checker, whose blocklist would otherwise have been the one
file in the repository carrying every string the repository forbids; its patterns
are stored encoded and decoded at import. Two tests do the same thing for the
same reason. This is a small point and it is a real one.}
\end{defect}

\begin{defect}{A traceability check that states its rule positively}
\symptom{The obvious design flags every numeric literal in prose and then excuses
most of them.}
\diagnosis{The exclusion list grows until the check means nothing. That is not a
hypothetical: it is the failure mode of essentially every lint rule that starts
by flagging everything.}
\resolution{State the rule positively. Three shapes count as a measurement: a
percentage, a probability shaped decimal, and a number sharing a line with a
measurement word. Years, dates, version strings, reference numbers, list markers,
code blocks and addresses are never flagged at all. A line is cleared either by a
result file reference or by an explicit annotation carrying a written reason, so
that waving the check through leaves a written trace with a name on it.}
\instead{Flag everything and maintain an exclusion list. Chosen against
deliberately, and the reasoning is in the checker's own docstring so that the
next person to be annoyed by a false negative knows what the trade was.}
\end{defect}

\begin{defect}{The traceability check passed on citations of files that did not
exist}
\symptom{Until the first result files existed the check could only look for the
\emph{shape} of a citation, and a check that matches a shape passes on a citation
of a file nobody committed.}
\diagnosis{That is the same thing as no citation at all, with more characters.}
\resolution{A resolve mode that walks the whole chain for every reference: the
path must exist, a manifest must sit beside it or above it, the manifest must
name a commit this repository can produce an object for, and the manifest must
not mark the run dirty. The scan and the resolve are separate so that the
scanner's own tests can run over strings with no repository behind them.}
\instead{Leaving it as a shape match. It had been green since the first phase and
would have stayed green forever, because it never asked the version control
system for anything.}
\end{defect}

\begin{defect}{The first genuinely red gate, and it was correct}
\symptom{The traceability job went red on a push, on a check that had passed
locally minutes earlier, reporting that a manifest's commit does not resolve in
this repository.}
\diagnosis{The job checked out at the default depth of one, so the clone held
only the tip commit and every citation of a run taken in an earlier commit
failed. The message was true of the clone and false of the repository.}
\resolution{Full history on that job, the same setting the ancestry job already
carried for the same reason: both checks ask questions about history, and history
is the input.}
\instead{Fixing it from the error message without reproducing it. The failure was
reproduced locally with a shallow clone before anything was changed, because a
failure that is fixed without being reproduced is a failure that is guessed at.

Worth recording for a second reason. It is the first time a gate in this project
went red on a real defect rather than on a deliberately broken branch, so the
argument for proving that each gate can fail now has a natural example beside its
manufactured ones.}
\end{defect}

\begin{defect}{The proof of failure exercise found two real problems}
\symptom{Each continuous integration job was deliberately broken on a scratch
branch, one at a time, to prove it can go red. Two of the runs found something
that was not the intended breakage.}
\diagnosis{The packaging job's first attempt assumed the installer would place
the console script at a fixed path, which is a guess and was wrong on the runner.
Separately, a deliberate taxonomy breakage failed both the reachability job and
the test job.}
\resolution{The packaging job now sets its own installer home and binary
directory, which both fixes the assumption and gives the built artefact a
genuinely clean environment to install into. The double failure was left alone,
because it is the correct outcome: the property is asserted in two independent
places on purpose, and a defect that only one of them noticed would mean the other
was not doing its job.}
\instead{Skipping the exercise on the grounds that all six jobs were green. A
gate that has only ever been green is indistinguishable from a gate that always
returns green, and the run links for each observed failure are recorded in the
repository so that the claim is checkable rather than asserted.}
\end{defect}

\begin{defect}{Placeholder modules that sank the coverage gate}
\symptom{Measured coverage of the library sat well below the gate at a point when
the library was a taxonomy, a small command line surface, and a set of empty
module placeholders.}
\diagnosis{The placeholders were written with a single future import each. That
is one executable statement per file, thirty of them, none ever imported.}
\resolution{Strip the import. The placeholders then have zero statements, which
is what the coverage gate should see: the phases that fill them bring their own
tests.}
\instead{Lowering the gate for the first phase, or excluding the placeholders in
the coverage configuration. Both would have needed to be undone later and one of
them would have been forgotten.}
\end{defect}

\begin{defect}{A reachability check that does more than it was asked to}
\symptom{Not a defect. A deliberate strengthening, recorded because it has a
sharp edge.}
\diagnosis{The rule that the taxonomy has exactly one definition in code is
enforceable mechanically by scanning the source and the scripts for label
literals written anywhere else.}
\resolution{It does that. The scan is limited to tokens that cannot be confused
with ordinary prose, meaning the hyphenated specification tokens and the
uppercase extras.}
\instead{Scanning for every token. Several specification tokens are also ordinary
English words, and a check that fires on one of them in a comment is a check
somebody turns off. A check that cries wolf is worse than no check, because it
trains its reader to ignore it.}
\end{defect}

\begin{defect}{A first remote that was abandoned}
\symptom{The repository was initially pushed to a remote whose name differed in
case and separator from the one intended.}
\diagnosis{A naming mistake, caught before any run link had been recorded
anywhere.}
\resolution{The remote was repointed before anything referenced it, so no
recorded continuous integration link refers to the abandoned repository.}
\instead{Deleting the abandoned remote. The credentials in use do not carry
delete scope, so it is left in place and recorded here rather than quietly
forgotten.}
\end{defect}

\chapter{What is still open}
\label{ch:still-open}

Everything in the previous chapters was fixed, refused, or recorded as a
deliberate decision with a cost. These are the ones that are still live, and they
are listed here rather than in an issue tracker because a defect record that
stops at the fixed ones is a marketing document.

\begin{defect}{The training manifest of the shipped model is marked dirty}
\symptom{The manifest of the training run that produced the currently shipped
model records that the working tree carried uncommitted changes.}
\diagnosis{The model was fitted in the same working session that authored the new
corpus templates. The templates were in the tree and not yet committed at the
moment of fitting, which is the natural order to work in and the wrong order to
record in.}
\resolution{None retroactively. Refitting now to obtain a clean manifest would be
regenerating an artefact to fix a flag, which is the same move the
reproducibility rules forbid for results, and it would produce a different model
from the one every committed evaluation result was measured against.}
\instead{Ignoring it, or refitting. What is done instead is to state it plainly in
three places: the main report's method chapter, its limitations chapter, and
here. What follows from it is bounded and worth being precise about. The model
artefacts are committed and are content addressed into every evaluation run
manifest, so \emph{which} model produced the headline numbers is not in doubt.
The corpus digest in the training manifest matches the digest in every evaluation
manifest, so the model was fitted on the same corpus it was evaluated against.
What is genuinely weaker is the development split evidence, which rests on a run
whose tree state was not pinned to a commit.

The procedural fix is in the future work list and it is one sentence: commit the
corpus change first, then train from the clean tree.}
\end{defect}

\begin{defect}{The evaluation runner still classifies every form twice}
\symptom{Described in full in the language model chapter. The excess disagreement
between a row's predicted label and its finding codes is on the order of two
fields in seven hundred for a non deterministic engine.}
\diagnosis{The runner classifies once directly and once through the audit engine,
and for a non deterministic engine those are separate draws.}
\resolution{Not applied. The fix is to pass the first pass's predictions into the
audit rather than letting the audit reclassify.}
\instead{Fixing it in the phase that measured it. Changing what the runner does
for every engine after seeing the results is a regeneration, and the honest place
for the fix is a phase that re-registers and re-runs.}
\end{defect}

\begin{defect}{A rule in the table that cannot fire}
\symptom{The Japanese postal mark rule, described in the normalisation chapter.}
\diagnosis{Normalisation treats a lone symbol character as a delimiter.}
\resolution{Not applied. Making it reachable means changing what counts as a
token character, which churns every golden snapshot and every committed fixture
expectation.}
\instead{A quiet fix inside an unrelated phase. It is a deliberate change with a
cost and it belongs to a phase that pays the cost on purpose.}
\end{defect}

\begin{defect}{The published package carries no model}
\symptom{An install from the published artefact runs the rule baseline, because
no model bundle travels with it.}
\diagnosis{Model binaries in a package index artefact are a size and licensing
decision that has not been made, and the tool is designed to work without one.}
\resolution{Documented behaviour rather than a silent fallback: the automatic
engine prefers the model when one loads and falls back to the rule baseline with
one printed line when none does.}
\instead{Failing when no model is present, or shipping a model without deciding
the question. The current state has a real cost and it is stated: the default
install is not the default engine that the measurements recommend, and closing
that gap is on the future work list.}
\end{defect}

\begin{defect}{Five cross repository issues, all open}
\symptom{Every issue filed against the evaluation harness is still open. One of
them, publication of that package to an index, is a hard blocker on tagging a
stable release here, because this project's own rule forbids a stable tag while
any dependency is a source reference rather than a released version.}
\diagnosis{The rule is a gate criterion rather than a preference, and it was
written before it became inconvenient.}
\resolution{Not applied. The ledger records each issue with its status and the
workaround, if any, that stands here in the meantime.}
\instead{Vendoring the three functions this project uses and removing the
dependency. That is the cheap move and it would destroy the evidence the boundary
exists to produce, which is the whole reason the dependency is external in the
first place. The rule holds even when holding it is the expensive option, which
is the only circumstance in which a rule tells you anything.}
\end{defect}

\begin{defect}{Nothing here has been measured on a real page}
\symptom{Every accuracy number in this project is a measurement on a corpus this
repository generated.}
\diagnosis{By design, because content from real pages never enters this
repository. The cost is that the strongest objection to every accuracy number is
that a generated corpus has a generator's habits, and the rule table's vocabulary
and the corpus's label text were written by the same person.}
\resolution{Not applied, and it is the single most valuable thing the project
could do next.}
\instead{Relaxing the synthetic data rule. The right shape is a small hand
labelled set of real pages with a published labelling protocol and a published
disagreement rate, which needs no scraping and no stored page content, and it
should arrive with its prediction registered first. The honest registration is
that the rule table's advantage will shrink.}
\end{defect}


\end{document}
